
\documentclass{article}
\usepackage{colortbl}
\usepackage{makecell}
\usepackage{multirow}
\usepackage{supertabular}

\begin{document}

\newcounter{utterance}

\centering \large Interaction Transcript for game `codenames', experiment `frequency\_high', episode 5 with reflexion{-}Qwen3.5{-}27B.
\vspace{24pt}

{ \footnotesize  \setcounter{utterance}{1}
\setlength{\tabcolsep}{0pt}
\begin{supertabular}{c@{$\;$}|p{.15\linewidth}@{}p{.15\linewidth}p{.15\linewidth}p{.15\linewidth}p{.15\linewidth}p{.15\linewidth}}
    \# & $\;$A & \multicolumn{4}{c}{Game Master} & $\;\:$B\\
    \hline

    \theutterance \stepcounter{utterance}  
    & & & \multicolumn{2}{p{0.3\linewidth}}{
        \cellcolor[rgb]{0.9,0.9,0.9}{
            \makecell[{{p{\linewidth}}}]{
                \texttt{\tiny{[GM$|$GM]}}
                {'hidden': {'team': ['life', 'moment', 'hospital', 'case'], 'innocent': ['eyes', 'dinner', 'store', 'luck'], 'opponent': [], 'assassin': ['voice']}, 'revealed': {'team': {'team': [], 'innocent': [], 'opponent': [], 'assassin': []}, 'opponent': {'team': [], 'innocent': [], 'opponent': [], 'assassin': []}}}
            }
        }
    }
    & & \\ \\

    \theutterance \stepcounter{utterance}  
    & & \multicolumn{4}{p{0.6\linewidth}}{
        \cellcolor[rgb]{0.9,0.9,0.9}{
            \makecell[{{p{\linewidth}}}]{
                \texttt{\tiny{[P1$\langle$GM]}}
                \texttt{You have attempted this same task before and did not succeed. These are the plans you wrote after those attempts.} \\
\texttt{Trial \#1: New plan: My previous strategy failed because I prioritized grouping multiple words under a single clue (e.g., "story" for life, moment, case), which increased the risk of my partner guessing a distractor or opponent word and ending the turn prematurely. I also failed to account for the specific semantic overlap between my targets and the remaining distractors (e.g., "heart" for "health"). Next time, I will adopt a conservative strategy: I will provide clues that relate to only one or two highly distinct target words to minimize ambiguity. I will explicitly cross{-}reference my potential clues against the remaining distractor and opponent lists to ensure no semantic overlap exists before submitting. For the final turn with "life, moment, hospital, case," I should have chosen a very specific clue like "nurse" (targeting only hospital) or "time" (targeting only moment) rather than a broad concept like "story," ensuring the partner guesses the intended word without triggering a turn{-}ending error.} \\
\\ 
\texttt{Trial \#2: New plan: My previous attempt failed because I only targeted two words ("life", "moment") with the clue "time", leaving "hospital" and "case" unguessed, which prevented me from achieving the goal of guessing all team words. Additionally, "time" is a broad concept that could ambiguously relate to distractors like "luck" (e.g., "time is luck") or "dinner" (e.g., "dinner time"), increasing the risk of a turn{-}ending error. Next time, I will prioritize a clue that covers all four target words or at least the remaining unguessed ones in a single turn to maximize efficiency. I will select a specific medical or legal term like "patient" or "file" that strongly links "hospital" and "case" while also being semantically distant from the distractors ("eyes", "dinner", "store", "luck") and the kill word ("voice"). If a single clue cannot safely cover all four, I will ensure the clue targets the most difficult words first, avoiding broad temporal concepts that invite ambiguity.} \\
\\ 
\\ 
\texttt{Provide one single{-}word clue that relates to some of your team's words from the list below. You can choose to how many and to which words your clue relates to.} \\
\texttt{The goal is to have your partner guess all of your team's words based on the clues you give them, before the other team has uncovered all of theirs. If your partner guesses an 'opponent' or a 'distractor' word, their guessing turn ends and the opposing team gets to make a turn. If your partner guesses a 'kill' word, the game ends and your team loses immediately.} \\
\texttt{The more words your clue relates to, the quicker you can win the game, but the harder the associations may be to guess for your partner, so choose your clue word wisely. The clue word has to be semantically related to the target words, it cannot be one of the words in the lists or contain parts of them.} \\
\\ 
\texttt{Always give your single{-}word clue and your comma{-}separated list of related target words in the following format and make your answers as short as possible, never include any other text than is required in this form:} \\
\\ 
\texttt{CLUE: <WORD>} \\
\texttt{TARGETS: <TARGETS>} \\
\\ 
\texttt{Your team words are: life, moment, hospital, case.} \\
\texttt{Your opponent's team words are: .} \\
\texttt{Distractor words are: eyes, dinner, store, luck.} \\
\texttt{The kill words are: voice.} \\
            }
        }
    }
    & \\ \\

    \theutterance \stepcounter{utterance}  
    & \multicolumn{4}{p{0.6\linewidth}}{
        \cellcolor[rgb]{0.9,0.9,0.9}{
            \makecell[{{p{\linewidth}}}]{
                \texttt{\tiny{[P1$\rangle$GM]}}
                \texttt{CLUE: patient} \\
\texttt{TARGETS: hospital, case} \\
            }
        }
    }
    & & \\ \\

    \theutterance \stepcounter{utterance}  
    & & & \multicolumn{2}{p{0.3\linewidth}}{
        \cellcolor[rgb]{0.9,0.9,0.9}{
            \makecell[{{p{\linewidth}}}]{
                \texttt{\tiny{[GM$|$GM]}}
                \texttt{patient} \\
            }
        }
    }
    & & \\ \\

    \theutterance \stepcounter{utterance}  
    & & & \multicolumn{2}{p{0.3\linewidth}}{
        \cellcolor[rgb]{0.9,0.9,0.9}{
            \makecell[{{p{\linewidth}}}]{
                \texttt{\tiny{[GM$|$GM]}}
                ['hospital', 'case']
            }
        }
    }
    & & \\ \\

    \theutterance \stepcounter{utterance}  
    & & & \multicolumn{2}{p{0.3\linewidth}}{
        \cellcolor[rgb]{0.9,0.9,0.9}{
            \makecell[{{p{\linewidth}}}]{
                \texttt{\tiny{[GM$|$GM]}}
                {'word': 'hospital', 'assignment': 'team'}
            }
        }
    }
    & & \\ \\

    \theutterance \stepcounter{utterance}  
    & & & \multicolumn{2}{p{0.3\linewidth}}{
        \cellcolor[rgb]{0.9,0.9,0.9}{
            \makecell[{{p{\linewidth}}}]{
                \texttt{\tiny{[GM$|$GM]}}
                {'word': 'case', 'assignment': 'team'}
            }
        }
    }
    & & \\ \\

    \theutterance \stepcounter{utterance}  
    & & & \multicolumn{2}{p{0.3\linewidth}}{
        \cellcolor[rgb]{0.9,0.9,0.9}{
            \makecell[{{p{\linewidth}}}]{
                \texttt{\tiny{[GM$|$GM]}}
                \texttt{opponent has won} \\
            }
        }
    }
    & & \\ \\

    \theutterance \stepcounter{utterance}  
    & \multicolumn{4}{p{0.6\linewidth}}{
        \cellcolor[rgb]{0.95,0.95,0.95}{
            \makecell[{{p{\linewidth}}}]{
                \texttt{\tiny{[P2$\rightarrow$P1]}}
                \texttt{Post{-}episode Reflexion (not sent to GM):} \\
\texttt{New plan: My previous attempt failed because I only targeted two words ("hospital", "case") with the clue "patient", leaving "life" and "moment" unguessed, which prevented me from achieving the goal of guessing all team words. Additionally, "patient" is too specific to the medical context and does not semantically bridge the gap to "life" or "moment". Next time, I will prioritize a single clue that semantically connects all four target words ("life", "moment", "hospital", "case") to maximize efficiency and ensure the turn is completed in one go. I will select a broad but safe concept like "event" or "story" (if not a distractor) or a specific medical/legal narrative term that encompasses the duration of a "life", a specific "moment", a "hospital" stay, and a "case" file, while strictly verifying that this clue does not trigger the distractors ("eyes", "dinner", "store", "luck") or the kill word ("voice"). If a single safe clue for all four is impossible, I will choose a clue that covers the three most distinct words, ensuring the remaining word is easily guessable in a subsequent turn without risking a turn{-}ending error.} \\
            }
        }
    }
    & & \\ \\

\end{supertabular}
}

\end{document}
